\chapter{随机过程、Markov、鞅、Poisson 与 Brownian}
\label{ch:stochastic-processes}

随机过程把“一个随机数”扩展为随时间变化的对象，但时间本身并没有告诉我们应该保留哪些信息。本章先区分有限维分布、样本路径和滤过，再用 Markov 性、鞅性、Poisson 到达和 Brownian 路径展示四种不同的压缩方式。贯穿全章的问题是：一个结论究竟依赖状态、信息集，还是依赖路径的正则性？

每个模型都同时配有一个可计算的基线和一个会失败的边界。这样读者在把过程用于价格、订单流或风险控制之前，能够先判断“模型结构”与“数据像不像”之间还差哪些假设。

\section{一个过程有三个不可互换的层次}

随机过程 $(X_t)_{t\in T}$ 是同一概率空间 $(\Omega,\mathcal F,\mathbb P)$ 上、取值于状态空间 $(E,\mathcal E)$ 的随机变量族。任取 $t_1<\cdots<t_k$，联合律
\[
  \mathcal L(X_{t_1},\ldots,X_{t_k})
\]
称为有限维分布；固定 $\omega$ 后，$t\mapsto X_t(\omega)$ 是样本路径。\textbf{选读：过程构造与版本。}有限维分布族若满足相容性，Kolmogorov 扩张定理可构造过程的概率律，但它并不自动给出连续路径。两个过程可以同分布却不是同一版本；若每个固定 $t$ 都几乎处处相等，称为版本，若整条路径除零测集外处处相等，才称不可区分。连续修改、右连续和 càdlàg 都是需要另证的路径正则性。

% Generated by tools/build_figures.py; edit figures/figure-specs.json instead.
\MFQFigure{tex/figures/generated/upper-ch11-filtration-timeline.pdf}{滤过把“何时知道什么”写进随机过程模型；信息集只能随时间扩张。}{upper-ch11-filtration-timeline}

% Short alias: % Short chapter-local alias; the cached asset remains registered in figures/figure-manifest.json.
% Generated by tools/build_figures.py; edit figures/figure-specs.json instead.
\MFQFigure{tex/figures/generated/upper-ch11-filtration-timeline.pdf}{滤过把“何时知道什么”写进随机过程模型；信息集只能随时间扩张。}{upper-ch11-filtration-timeline}



滤过 $(\mathcal F_t)_{t\ge0}$ 满足 $\mathcal F_s\subseteq\mathcal F_t$（$s\le t$），表示信息增长。自然滤过是 $\sigma(X_s:s\le t)$；实际定理常再完备化并取右连续修改。$X_t$ 对 $\mathcal F_t$ 可测称为适应。交易指令若必须在 $t$ 前确定，常需要可预测性而不只是适应性。随机时刻 $\tau$ 若 $\{\tau\le t\}\in\mathcal F_t$，才是停止时刻；“全样本最优离场点”通常不是。


\begin{example}[信息层次]
同一条价格序列既可被当作有限维向量做协方差估计，也可被当作路径讨论最大回撤。前者不能替代后者。若用收盘后才知道的日最高价决定盘中下单，变量虽属于终端 $\mathcal F_T$，却不属于当时的信息集。
\end{example}

“同一个过程”因此至少有三个问题：有限个时点的联合律是什么，单条路径具有什么正则性，以及在每个决策时点允许使用哪些信息。若只验证边缘分布，可能错过依赖结构；若只画路径，可能把一次样本的形状误当成总体性质；若只写终端变量，又会把未来信息偷偷带入策略。后面每个模型都要先回答这三个问题，再谈参数估计。

\section{Markov 状态压缩与鞅公平性}

齐次离散 Markov 链满足：对任意使条件事件概率为正的历史
$X_0=i_0,\ldots,X_{n-1}=i_{n-1},X_n=i$，
\[
 \mathbb P(X_{n+1}=j\mid
 X_0=i_0,\ldots,X_{n-1}=i_{n-1},X_n=i)
 =\mathbb P(X_{n+1}=j\mid X_n=i)=P_{ij}.
\]
转移矩阵 $P$ 每行非负且和为一，Chapman--Kolmogorov 关系给出 $P^{m+n}=P^mP^n$。Markov 性不是“相邻独立”，而是当前状态已包含预测未来所需的历史摘要。若摘要遗漏库存、剩余期限或波动状态，扩充状态往往比强行接受 Markov 假设更诚实。

\begin{proposition}[有限状态链的矩阵递推]
若 $\mu_n$ 是时刻 $n$ 的行分布，齐次转移矩阵为 $P$，则
\[
  \mu_{n+k}=\mu_nP^k.
\]
若 $\pi P=\pi$，则从 $\pi$ 出发的过程在每个时点都具有同一分布；若链有限、不可约且非周期，则存在唯一 $\pi$，并且任意初始分布都满足 $\mu_n\to\pi$。
\end{proposition}
\begin{proof}
条件概率的全概率分解给出 $\mu_{n+1}(j)=\sum_i\mu_n(i)P_{ij}$，即 $\mu_{n+1}=\mu_nP$。归纳得到矩阵幂公式；稳态结论直接由 $\pi P=\pi$ 得到。有限不可约非周期链的收敛来自 Perron--Frobenius 定理：特征值 1 是单的，其余特征值模小于 1。周期链会保留单位圆上的其他特征值，可约链则可能有多个闭沟通类，这正是后面两个反例的来源。
\end{proof}

实验取
\[
P=\begin{pmatrix}0.8&0.2\\0.3&0.7\end{pmatrix},\qquad \mu_0=(1,0).
\]
令 $p_n$ 为时刻 $n$ 处于第一状态的概率。按上一期状态分解，
\[
p_{n+1}=0.8p_n+0.3(1-p_n)=0.3+0.5p_n.
\]
稳态解满足 $p=0.3+0.5p$，故 $p=0.6$。减去稳态得到 $p_{n+1}-0.6=0.5(p_n-0.6)$，逐次代入、利用 $p_0=1$，
\[
p_n=0.6+0.4\cdot0.5^n,\quad p_2=0.7,\quad p_5=0.6125.
\]
因此 $\pi=(0.6,0.4)$，且 $\mu_0P^n-\pi=(0.4,-0.4)0.5^n$。


这条显式公式把“趋于稳态”变成可计算的误差界。若 $i\to j$ 与 $j\to i$ 都有正概率路径，称两状态沟通；沟通关系把状态分成沟通类。状态若以概率 1 返回称常返，否则称暂留；常返状态的期望返回时间有限时称正常返。有限不可约链中所有状态正常返且稳态唯一，若再非周期，才有 $\mu P^n\to\pi$。若 $\pi_iP_{ij}=\pi_jP_{ji}$ 对所有 $i,j$ 成立，称详细平衡；它蕴含稳态但不是必要条件。本例双向流均为 $0.6\times0.2=0.4\times0.3$，详细平衡流量为 $0.120000$。

\noindent\textbf{两个不能省略的反例。}\quad
交替链 $P=\bigl(\begin{smallmatrix}0&1\\1&0\end{smallmatrix}\bigr)$ 有稳态 $(1/2,1/2)$，但从状态 0 出发的分布在 $(1,0)$ 与 $(0,1)$ 间振荡，周期为 2。可约链 $P=\bigl(\begin{smallmatrix}1&0\\1/2&1/2\end{smallmatrix}\bigr)$ 有吸收类；从状态 1 最终吸收于 0 的概率为 1，期望等待时间为 2。故“算出稳态”不等于“任意初值收敛”。

命中概率和期望命中时长来自边界值问题。公平随机游走在 $\{0,1,\ldots,N\}$ 上，从 $i$ 出发，令 $h_i=\mathbb P_i(\tau_N<\tau_0)$、$m_i=\mathbb E_i(\tau_0\wedge\tau_N)$，则
\[
h_i=\tfrac12h_{i-1}+\tfrac12h_{i+1},\quad h_0=0,h_N=1,
\]
\[
m_i=1+\tfrac12m_{i-1}+\tfrac12m_{i+1},\quad m_0=m_N=0.
\]
二阶差分方程给出 $h_i=i/N$、$m_i=i(N-i)$。取 $N=4,i=2$，赌徒破产命中概率与期望时长为 $0.500000$ 与 $4.000000$。

最后，令独立 Bernoulli$(1/2)$ 变量 $U,V$ 生成 $(X_0,X_1,X_2)=(U,V,U)$。给定 $X_1$ 时与再给定 $X_0=1$ 时，$X_2=1$ 的条件概率分别为 $0.5$ 与 $1$。因此相邻独立不推出 Markov 性。

\MFQLead{鞅性依赖信息集与可积性}

可积、适应过程 $(M_n)$ 若满足 $\mathbb E(M_{n+1}\mid\mathcal F_n)=M_n$，称为鞅；将等号改为 $\ge$ 或 $\le$ 得次鞅或上鞅。Doob 分解说明合适的次鞅可写为“鞅 + 可预测递增补偿项”。对平方可积鞅，$\langle M\rangle_n=\sum_{k=1}^n\mathbb E[(\Delta M_k)^2\mid\mathcal F_{k-1}]$ 累积下一步条件方差；它与从已实现路径计算的 $[M]_n=\sum(\Delta M_k)^2$ 不能混称。

对称随机游走 $S_n=\sum_{k=1}^n\xi_k$ 满足 $\mathbb E(S_{n+1}\mid\mathcal F_n)=S_n$。又因 $\mathbb E(\xi_{n+1}\mid\mathcal F_n)=0$、$\mathbb E(\xi_{n+1}^2\mid\mathcal F_n)=1$，
\[
\mathbb E(S_{n+1}^2-(n+1)\mid\mathcal F_n)=S_n^2-n,
\]
故 $S_n^2-n$ 是鞅，而 $S_n^2$ 是次鞅；两步随机游走有 $\langle S\rangle_2=2$。类似地，Poisson 过程的可预测补偿子为 $\lambda t$，所以 $M_t=N_t-\lambda t$ 是鞅，且 $\langle M\rangle_t=\lambda t$。取 $\lambda t=6$，Poisson 补偿鞅的可预测二次变差为 $6.000000$。

可选停止定理不是“任何下注策略都保期望”。有界停止时刻足够；其他版本需要一致可积、受控增量或可积停止时间等条件。加倍下注在首次胜利时赚 1，但若截断于 $m$ 轮，未胜路径损失 $2^m-1$，概率 $2^{-m}$，于是
\[
\mathbb E W_m=(1-2^{-m})-2^{-m}(2^m-1)=0.
\]
虽然 $W_m\to1$ 几乎处处，尾部负损失对期望的贡献为 $1-2^{-m}\to1$，不能交换极限与期望；这正是缺少一致可积性的代价。

若 $\tau$ 是有界停止时刻，鞅 $(M_n)$ 的停止过程
$M_{n\wedge\tau}$ 仍是鞅，因此
\[
  \mathbb E[M_\tau]=\mathbb E[M_0].
\]
证明只用塔式性质：对每一步把新增项写成
$\mathbf1_{\{\tau\ge k\}}(M_k-M_{k-1})$，指标属于 $\mathcal F_{k-1}$，条件期望为零；有界性允许有限次求和。若 $\tau$ 不有界，必须再证明一致可积或足够强的尾部控制，才能把 $n\wedge\tau$ 的极限移入期望。这个推导正好解释加倍下注的漏洞：几乎处处收敛并不提供所需的可积尾部界。

\section{Poisson：基线、变换与聚集失效}

强度 $\lambda$ 的齐次 Poisson 过程从零开始，具有独立平稳增量，且 $N_t-N_s\sim\mathrm{Poisson}(\lambda(t-s))$。因此 $\mathbb E N_t=\operatorname{Var}N_t=\lambda t$，相邻到达等待时间独立且服从 $\mathrm{Exp}(\lambda)$。取 $\lambda=2,t=3$，模拟计数的均值与方差为 $5.980833$ 与 $5.988466$，解析目标均为 6。

三条变换应能直接推导：独立强度 $\lambda_1,\lambda_2$ 的过程叠加后强度相加；对每次到达以概率 $q$ 独立保留，保留与拒绝流分别是强度 $q\lambda$、$(1-q)\lambda$ 的独立 Poisson 过程；非齐次强度 $\lambda(t)$ 的累计强度为 $\Lambda(t)=\int_0^t\lambda(s)\,ds$，经时间变换可还原单位速率过程。独立基准取等待率 2、叠加率 $2+3$、基准率 5 且 $q=0.3$，以及 $\lambda(t)=2+t$、$t\le3$，得到均值/方差 $0.5,0.25$，叠加强度 5，稀疏化强度 $1.5,3.5$，累计强度 $10.5$。

从独立小区间增量还可以得到观测窗口 $[0,T]$ 内的似然。若观察到 $N_T=n$ 次到达，则
\[
 \mathbb P(N_T=n)=e^{-\lambda T}\frac{(\lambda T)^n}{n!},
 \qquad
 \ell(\lambda)=n\log\lambda-\lambda T+\text{const}.
\]
当 $n>0$ 时，令导数为零得到 $\widehat\lambda=n/T$，二阶导数为 $-n/\lambda^2<0$；$n=0$ 时最大值在允许的边界 $\lambda=0$ 取得。因此 Poisson 强度估计不是只能从模拟均值获得；它有明确的统计目标和不确定性。若事件被合并、漏记或存在日内强度，直接使用 $n/T$ 仍可能是对错误观测过程的精确估计。

若订单计数方差显著大于均值、等待时间非指数、增量自相关或日内强度变化，齐次 Poisson 基线失效。非齐次 Poisson 可处理季节性，却仍不能产生自激聚集；Hawkes、Cox 或状态切换模型才是后续候选。模型升级应由残差诊断和样本外预测推动，而非因为模型名字更复杂。

\section{Brownian：高斯结构、命中与粗糙路径}

标准 Brownian 运动从零开始，路径几乎处处连续，且不相交区间上的增量独立并满足
\[
W_t-W_s\sim\mathcal N(0,t-s),\qquad \operatorname{Cov}(W_s,W_t)=\min(s,t).
\]
它具有缩放性 $(W_{ct})_{t\ge0}\overset d=(\sqrt c\,W_t)_{t\ge0}$；取 $c=4,t=1/4$，两边方差都为 1。\par\medskip\noindent\textbf{选读：随机游走的函数极限。}本段只说明 Brownian 运动为何能作为缩放随机游走的极限；完整证明需函数空间弱收敛理论，不作为本章基础练习要求。Donsker 不变性原理说，若 $(\xi_k)$ 独立同分布、均值为零且方差为 1，令 $S_m=\sum_{k=1}^m\xi_k$，则折线插值
\[
 W_n(t)=\frac{S_{\lfloor nt\rfloor}
 +(nt-\lfloor nt\rfloor)\xi_{\lfloor nt\rfloor+1}}{\sqrt n},
 \qquad 0\le t\le1,
\]
在 $C[0,1]$ 的弱收敛意义下趋于 Brownian 运动。若改用
$S_{\lfloor nt\rfloor}/\sqrt n$，得到的是 càdlàg 阶梯过程，应在
$D[0,1]$ 中表述。两种版本都不是有限 $n$ 的逐路径相等，也不覆盖无限方差或强依赖增量。对 $n=100$ 的对称增量，归一化终值方差为 1，Donsker 四阶矩基线为 $2.980000$。在 $t=(0.25,0.5,1)$ 上，经验协方差与 $\min(s,t)$ 的最大差为 $0.006948$。

反射原理给出（$a>0$）
\[
\mathbb P\!\left(\sup_{0\le s\le t}W_s\ge a\right)
=2\mathbb P(W_t\ge a)=2\{1-\Phi(a/\sqrt t)\}.
\]
当 $a=t=1$ 时，反射原理给出命中概率 $0.317311$。这计算的是连续监控命中概率；用日收盘离散观察会漏掉区间内越界。

反射原理的关键是把第一次越过 $a$ 的路径在命中时刻之后关于水平 $a$ 反射；对已经命中 $a$ 而终点低于 $a$ 的路径，反射后的终点为 $2a-W_t>a$；这与终点高于 $a$ 的路径一一对应，并由强 Markov 性和增量对称性保持概率。于是命中事件被拆成两个同概率尾部，而不是凭经验把终点尾概率乘二。这个证明也说明离散观察的误差来自“在网格之间越过后又回到阈值内”的路径，而不是来自正态尾概率计算本身。

Brownian 路径几乎处处不可微且总变差无限，但沿网格的二次变差收敛到区间长度。等距 $n$ 段上，绝对增量和的期望为 $n\mathbb E|N(0,1/n)|=\sqrt{2n/\pi}$；当 $n$ 从 16 增至 256，总变差基线从 $3.191538$ 增至 $12.766153$。相反，256 段二次变差模拟均值为 $1.000731$、方差为 $0.007781$，解析方差为 $2/256$。有限总变差过程与 Brownian 的这种差异，是 Itô 积分不能照搬普通 Riemann--Stieltjes 微积分的根源。

对分割 $\Pi_n=\{0=t_0<\cdots<t_n=t\}$，二次变差和为
\[
 Q_n=\sum_{k=1}^n(W_{t_k}-W_{t_{k-1}})^2.
\]
独立增量给出 $\mathbb E Q_n=t$；当网格最大长度趋于零时，方差
$2\sum_k(t_k-t_{k-1})^2$ 也趋于零，所以 $Q_n\to t$ 在 $L^2$ 中成立。这个可计算的结论比“路径很粗糙”更精确，也正是连续时间风险模型需要的积分基础。

\begin{note}
Brownian、Poisson 与 Markov 都是条件结构，不是资产收益的事实标签。量化研究中必须写清状态、滤过、观测频率、参数稳定性、交易时钟及模型失败后的替代方案。
\end{note}

\section{命中方程如何解出来}

对公平游走，$h_{i+1}-h_i=h_i-h_{i-1}$，因此相邻差为常数 $c$。由 $h_0=0$ 得 $h_i=ci$，由 $h_N=1$ 得 $c=1/N$。期望时长方程则给 $m_{i+1}-2m_i+m_{i-1}=-2$；因为 $i^2$ 的二阶差分为 2，试取 $m_i=-i^2+ai+b$，代入两端边界得到 $b=0,a=N$。这样得到 $i(N-i)$，而不是只靠记住答案。

在有限吸收链里，还需确认期望有限。每 $N$ 步至少有 $2^{-N}$ 的概率连续向某一边移动并吸收，因此存活块数被几何尾界控制，期望有限。notebook 将直接解有限线性方程并模拟命中时长；靠近边界的起点更快退出，中间起点最慢。

\noindent\textbf{延伸阅读。}Markov 链参见\cite{norris1997markov}；函数空间弱收敛与 Donsker 定理参见\cite{billingsley1999convergence}。

\section{笔面试与研究训练}
\begingroup\small
\subsection*{口述概念（3 题）}
\begin{enumerate}[nosep,leftmargin=*]
  \item 区分有限维分布、版本、不可区分和连续修改；为何边缘分布相同不够？
  \item 区分适应、可预测与停止时刻，并给出一个未来函数的交易例子。
  \item 稳态分布、极限分布、不可约、周期和正常返之间是什么关系？
\end{enumerate}
\subsection*{笔试推导（3 题）}
\begin{enumerate}[nosep,leftmargin=*,start=4]
  \item 解公平随机游走在 $\{0,\ldots,N\}$ 上先到 $N$ 的概率和吸收期望时长。
  \item 证明 $S_n^2-n$ 是鞅，并解释加倍下注为何不能用可选停止推出稳赚。
  \item 用反射原理推导 $\mathbb P(\sup_{s\le t}W_s\ge a)$，再说明离散监控为何低估越界概率。
\end{enumerate}
\subsection*{数值编程（3 题）}
\begin{enumerate}[nosep,leftmargin=*,start=7]
  \item 对两状态链核对 $P^5$ 与稳态，并构造周期链展示“有稳态但不收敛”。
  \item 独立核对 Poisson 的等待、叠加、稀疏化和非齐次累计强度账本。
  \item 模拟 Brownian 协方差、二次变差与绝对增量和；分别写出解析 oracle 和容差来源。
\end{enumerate}
\subsection*{研究判断（3 题）}
\begin{enumerate}[nosep,leftmargin=*,start=10]
  \item 订单到达呈日内季节性和聚集性时，如何逐级检验并替换齐次 Poisson？
  \item 某收益状态模型训练集拟合良好，却在制度切换后失效。如何判断是转移概率漂移还是状态定义不充分？
  \item 连续 Brownian 风险模型用于离散、跳跃、含交易停牌的市场时，应报告哪些模型风险？
\end{enumerate}
\endgroup

\section{分级提示}
\begingroup\small
\begin{enumerate}[nosep,leftmargin=*]
  \item 固定 $t$ 的零测集与对所有 $t$ 同时成立的零测集并不相同。
  \item 问“决策形成时已知什么”，而不只问变量最终是否可测。
  \item 先画沟通类；用交替链拆开稳态与收敛；详细平衡检查双向概率流。
  \item 为 $h_i,m_i$ 写一步递推和两个吸收边界，解一阶与二阶差分。
  \item 展开平方并条件化；Poisson 减去累计强度；加倍策略检查一致可积性。
  \item 把首次越过后的路径围绕水平 $a$ 反射；离散网格看不到两点之间的最大值。
  \item 同时打印偶数步与奇数步分布，不要只打印求出的 $\pi$。
  \item 等待时间用 $1/\lambda,1/\lambda^2$；叠加强度相加；稀疏化强度乘保留概率。
  \item 缩放比较方差；Donsker 终值四阶矩为 $3-2/n$；其余目标分别为 $\min(s,t)$、1 与 $\sqrt{2n/\pi}$。
  \item 依次检查分时强度、过度离散、等待时间、增量相关和样本外残差。
  \item 比较滚动转移矩阵与条件在扩充状态后的残差；不要只增加隐藏状态数量。
  \item 报告跳跃、离散监控、厚尾、停牌时钟、参数漂移和对冲/执行误差。
\end{enumerate}
\endgroup


\noindent\textbf{本章要点。}有限维分布、路径和信息集回答不同问题；Markov 链把历史压缩进状态，鞅把公平性绑定到滤过与可积性，Poisson 与 Brownian 分别提供跳跃和连续噪声基线。任何动态结论都应同时写出状态、信息集、增量结构和失败边界，而不能只凭一张路径图给模型贴标签。
